\scanpage{18}
\mspage{3}
``wholly degenerate'' generalized boundary). When we throw in the rigidifications and go over to isomorphism classes, this should give the desired retraction. However, the geometric situation is a lot richer still, as the compact surface without boundary obtained through surgery is endowed, not only with a system of $\nu$ points, but moreover with a system of mutually disjoint \uline{discs} around these points. The shape of these discs is by no means arbitrary - we'll say that a system of discs around $\nu$ points on a compact \inserted{conformal} surface $X^{\wedge}$ without boundary is ``admissible'', if the situation can be obtained as above (up to isomorphism) from surgery, starting with a compact conformal surface $X$ with boundary. (NB Among the given ``discs'', we should allow that some should be reduced to their center - we'll call them ``degenerate''.) The condition of admissibility can be expressed intrinsically, by stating that \deleted{the} for every non-degenerate component $\Gamma_i$ of the system of boundaries of those discs, the two operations we got of the standard circle group (of complex numbers of module $1$) upon $\Gamma_i$, by using the fact that it is (on the one hand) the boundary of the disc $D_i$, and (on the other hand) that it is a component of the boundary of the hyperbolic surface $X^{\wedge}\backslash(\bigcup_j D_j^{o})$, should be the same. When $X^{\wedge}$ and the points $s_i$ on $X^{\wedge}$ are given, the possible \inserted{admissible} systems of discs around the points $s_i$ depend on $\nu$ parameters - and the first idea which flips to mind to give a more precise meaning to these ``parameters'', is to view them as being the ``radii'' of those discs. But then we'll have to define what we mean by these !
\margin{!}

The idea here is that, when we have a conformal disc $D$ and an \unsure{interior} point $s$ of $D$, then $D$ may be viewed as canonically embedded in \unsure{the} tangent space $T_s$ to $D$ at $s$, as the ``unit disc'' at $s$. Thus, in the situation above of an\deleted{\unsure{just}} admissible system of discs $(D_i)_{i\in I}$ around $(s_i)_{i\in I}$, for every $s_i$ corresponding to a non-degenerate $D_i$, we get a canonical disc
\[
\Delta_i\subset T_{s_i}
\]
in the tangent space - and of course, for degenerate $D_i$, we'll take $\Delta_i$ to be degenerate too. The discs we get in a\deleted{\unsure{y}} given $T_{s_i}$ (for a fixed system $(s_i)$, and a variable \inserted{admissible} system of discs around these $s_i$) are all discs in the strict euclidean sense, given by an unequality
\[
|z|\le r_i\ ,
\]
where $z\mapsto|z|$ denotes some hermitian metric on $T_{s_i}$ compatible with the conformal structure - this metric being \deleted{defined}\inserted{unique} up to a scalar factor. the set\inserted{$R_i$} of all those possible discs (the non-degenerate ones\deleted{\unsure{,}} say) may be viewed in a natural way as a ``torsor'' (= principal homogeneous space)
% archive foliation: 17
